\chapter{Reaction Theory I}
\marginnoteinline{Source date: 01/07/2010.  The handwritten references appear to include Foderaro, Chapter 4, and a second chapter reference that is not fully legible.}

\sourcepage{1}
\section{Goal and scattering geometry}
Beginning with partial-wave analysis, review several simple examples of computing cross sections for elastic scattering.  Thereafter, derive the single-level Breit--Wigner formula for a resonance.

The sketch shows a plane incident beam, a target potential of characteristic radius $10^{-8}$--$10^{-12}\unit{cm}$, and a detector at a macroscopic radius of order $100\unit{cm}$.  The incident wave is
\begin{equation*}
  \psi_{\mathrm{inc}}\sim \ee^{\ii kz}=\ee^{\ii kr\cos\theta}.
  \tag{1}
\end{equation*}
At the detector, the scattered wave has the form
\begin{equation*}
  \psi_{\mathrm{sc}}\sim f(\theta)\frac{\ee^{\ii kr}}{r}.
  \tag{2}
\end{equation*}
Azimuthal symmetry is assumed, so $f$ depends only on $\theta$.  The factor $r^{-1}$ conserves radial flux.

\sourcepage{2}
\section{Differential cross section from current balance}
The differential scattering cross section is the probability that an incident particle scatters through the solid-angle element $\dd\Omega$.  Particle balance gives
\begin{equation*}
  J_{\mathrm{inc}}\dd\sigma=J_{\mathrm{sc}}\dd A.
  \tag{3}
\end{equation*}
Since
\begin{equation*}
  \dd A=r^2\sin\theta\dd\theta\dd\phi=r^2\dd\Omega,
  \tag{4}
\end{equation*}
we have
\[
  J_{\mathrm{inc}}\dd\sigma=J_{\mathrm{sc}}r^2\dd\Omega
\]
and therefore
\begin{equation*}
  \frac{\dd\sigma}{\dd\Omega}
  =\frac{J_{\mathrm{sc}}r^2}{J_{\mathrm{inc}}}.
  \tag{5}
\end{equation*}

Recall
\[
  J=\frac{\hbar}{2m\ii}(\psi^*\psi'-\psi^{*\prime}\psi).
\]
For the incident plane wave,
\[
  J_{\mathrm{inc}}=\frac{\hbar k}{m}=v.
\]
For the outgoing spherical wave $f(\theta)\ee^{\ii kr}/r$, the radial-current calculation gives
\begin{equation*}
  J_{\mathrm{sc}}=v\frac{|f(\theta)|^2}{r^2}.
  \tag{6}
\end{equation*}
Hence
\begin{equation*}
  \frac{\dd\sigma}{\dd\Omega}=|f(\theta)|^2.
  \tag{7}
\end{equation*}

\sourcepage{3}
From a previous result,
\begin{equation*}
  f(\theta)=\frac{1}{k}\sum_{\ell=0}^{\infty}(2\ell+1)
  \ee^{\ii\delta_\ell}\sin\delta_\ell
  P_\ell(\cos\theta).
  \tag{7a}
\end{equation*}
Therefore
\begin{equation*}
  \sigma(\theta)
  =\frac{1}{k^2}\left|
  \sum_{\ell=0}^{\infty}(2\ell+1)
  \ee^{\ii\delta_\ell}\sin\delta_\ell
  P_\ell(\cos\theta)
  \right|^2.
  \tag{8}
\end{equation*}
The total elastic scattering cross section is
\[
  \sigma_s=\int_0^\pi 2\pi\sin\theta\dd\theta\,\sigma(\theta).
\]
Using Legendre orthogonality,
\begin{equation*}
  \sigma_s=\frac{4\pi}{k^2}
  \sum_{\ell=0}^{\infty}(2\ell+1)\sin^2\delta_\ell.
  \tag{9}
\end{equation*}

In general, expand
\[
  f(\theta)=\sum_{\ell=0}^{\infty}(2\ell+1)f_\ell(k)P_\ell(\cos\theta),
\]
where $f_\ell$ is the partial-wave amplitude.  From above,
\begin{equation*}
  f_\ell(k)=\frac{1}{k}\ee^{\ii\delta_\ell}\sin\delta_\ell
  =\frac{1}{2\ii k}\left[S_\ell(k)-1\right],
  \qquad
  S_\ell(k)=\ee^{2\ii\delta_\ell}.
  \tag{10}
\end{equation*}

\sourcepage{4}
\section{Infinite repulsive sphere: setup}
Consider the infinite repulsive sphere or ``low neutron potential''
\[
  V(r)=
  \begin{cases}
    \infty, & r\le b,\\
    0, & r>b.
  \end{cases}
\]
The inner wave function is identically zero.  The outer wave function is a linear combination of $j_\ell(kr)$ and $n_\ell(kr)$.  Continuity at the surface requires
\[
  \psi_\ell^{\mathrm{out}}(b)=0.
\]
Far from the origin, the outer wave must have the phase-shifted form
\begin{equation*}
  \psi_\ell^{\mathrm{out}}
  =\frac{1}{kr}\sin\left(kr-\frac{\ell\pi}{2}+\delta_\ell\right).
  \tag{11}
\end{equation*}
The large-$r$ asymptotic forms are
\begin{align*}
  j_\ell(kr)&\sim \frac{1}{kr}\cos\left(kr-\frac{(\ell+1)\pi}{2}\right),\\
  n_\ell(kr)&\sim \frac{1}{kr}\sin\left(kr-\frac{(\ell+1)\pi}{2}\right).
  \tag{12}
\end{align*}

\sourcepage{5}
Expanding the shifted sine in exponentials and equating the incoming and outgoing exponential coefficients with a general combination
\[
  a_\ell j_\ell(kr)+b_\ell n_\ell(kr)
\]
gives
\begin{equation*}
  b_\ell=\sin\delta_\ell,
  \tag{15}
\end{equation*}
\begin{equation*}
  a_\ell=\cos\delta_\ell.
  \tag{16}
\end{equation*}

\sourcepage{6}
Thus the general exterior wave function is
\begin{equation*}
  \psi_\ell^{\mathrm{out}}(r)
  =\cos\delta_\ell\,j_\ell(kr)
  -\sin\delta_\ell\,n_\ell(kr).
  \tag{17}
\end{equation*}
For a finite interior potential, continuity of $\psi$ and $\psi'$ at $r=b$ gives
\begin{align*}
  \psi_\ell^{\mathrm{in}}(b)
  &=\cos\delta_\ell\,j_\ell(kb)
  -\sin\delta_\ell\,n_\ell(kb),\\
  \psi_\ell^{\mathrm{in}\prime}(b)
  &=\cos\delta_\ell\,j_\ell'(kb)
  -\sin\delta_\ell\,n_\ell'(kb).
\end{align*}
Define the logarithmic derivative
\begin{equation*}
  D\equiv\left[\frac{\dd\psi_\ell^{\mathrm{in}}/\dd r}
  {\psi_\ell^{\mathrm{in}}}\right]_{r=b}.
  \tag{18}
\end{equation*}
Then
\begin{equation*}
  D=\frac{\cos\delta_\ell\,j_\ell'(kb)-\sin\delta_\ell\,n_\ell'(kb)}
  {\cos\delta_\ell\,j_\ell(kb)-\sin\delta_\ell\,n_\ell(kb)}.
  \tag{19}
\end{equation*}
Solving for the phase shift,
\begin{equation*}
  \tan\delta_\ell
  =\frac{j_\ell'(kb)-D j_\ell(kb)}
  {n_\ell'(kb)-D n_\ell(kb)}.
  \tag{20}
\end{equation*}
The notes emphasize that $D$ still has to be computed or otherwise known.

\sourcepage{7}
For the solid sphere, the interior radial wave function is zero everywhere.  Therefore
\[
  \psi_\ell^{\mathrm{out}}(kb)=0
  =\cos\delta_\ell\,j_\ell(kb)-\sin\delta_\ell\,n_\ell(kb),
\]
and
\begin{equation*}
  \tan\delta_\ell=\frac{j_\ell(kb)}{n_\ell(kb)}.
  \tag{21}\label{eq:rt1-hard-sphere-phase}
\end{equation*}

For low energies, $kb\ll1$, and only the $\ell=0$ ($s$-wave) term matters.  The small-argument forms are
\begin{align*}
  j_\ell(x)&\sim\frac{x^\ell}{1\cdot3\cdot5\cdots(2\ell+1)},\\
  n_\ell(x)&\sim-\frac{1\cdot3\cdot5\cdots(2\ell-1)}{x^{\ell+1}}.
  \tag{22}
\end{align*}
Hence
\begin{equation*}
  \tan\delta_\ell
  \sim-\frac{(kb)^{2\ell+1}}
  {\left[1\cdot3\cdot5\cdots(2\ell-1)\right]^2(2\ell+1)}.
  \tag{23}
\end{equation*}
For $kb\ll1$, the $\ell=0$ term dominates:
\[
  \tan\delta_0\simeq-kb,
  \qquad
  \delta_0\simeq-kb.
\]
Thus
\begin{equation*}
  \sigma(\theta)=\frac{1}{k^2}\sin^2\delta_0\simeq b^2,
  \tag{24}
\end{equation*}
and
\begin{equation*}
  \sigma_s=4\pi b^2,
  \tag{25}
\end{equation*}
four times the classical geometric cross section.

\sourcepage{8}
For high energies, use Eq.~\eqref{eq:rt1-hard-sphere-phase} to write
\[
  \sin\delta_\ell
  =\frac{j_\ell(kb)}{\sqrt{j_\ell^2(kb)+n_\ell^2(kb)}},
\]
so
\begin{equation*}
  \sin^2\delta_\ell
  =\frac{j_\ell^2(kb)}{j_\ell^2(kb)+n_\ell^2(kb)}.
  \tag{26}
\end{equation*}
Using the large-$kb$ asymptotic forms,
\begin{equation*}
  \sin^2\delta_\ell
  \simeq\cos^2\left(kb-\frac{(\ell+1)\pi}{2}\right).
  \tag{27}
\end{equation*}
The partial-wave sum is cut off semiclassically at $\ell\sim kb$.  Averaging the rapid oscillation, $\cos^2\to1/2$, gives
\begin{align*}
  \sigma_s
  &\simeq\frac{4\pi}{k^2}
  \sum_{\ell=0}^{kb}(2\ell+1)
  \cos^2\left(kb-\frac{(\ell+1)\pi}{2}\right)\\
  &\simeq\frac{4\pi}{k^2}
  \int_0^{kb}(2\ell+1)\frac12\dd\ell\\
  &\simeq 2\pi b^2.
\end{align*}
Thus the high-energy quantum result approaches twice the classical geometric cross section.
